\section{Risks and required behaviour}
The public, staff, and restricted-access zones organise this model; they are not a formal classification of aviation security restricted areas. A single fire compartment SP0 and Polish occupancy category ZL III are assumed: a low-rise building with no overnight accommodation and no room for more than 50 people who are not regular users. IT is an ordinary locked room, without a separate fire compartment or extinguishing system.

\begin{tabularx}{\linewidth}{@{}P{32mm}Y Y@{}}
\toprule Event & Protection and limits & Failure response\\\midrule
Unauthorised access & Card, PIN for SOC/IT, door state, and associated image. A card does not prove who physically passed through. & Local permissions; an offline controller needs confirmed revocation updates or physical supervision.\\
Intrusion or hold-up & Two IDS partitions, 2EOL circuits, and a hold-up button; no complete fence detection system. & Local alarm and keypad; patrol replaces unavailable shared monitoring.\\
Fire & Detection and notification in every room; no delay for video verification. & Dedicated panel battery and mechanical exits; LAN and AI are outside the safety path.\\
IT or power loss & UPS, local ACS permissions, and camera SD recording while PoE remains available. & Loss of PoE stops video and electronic ACS; free egress remains available.\\
False analytics alarm & One rule in one camera, assessed using TP/FP/FN and operator workload. & Disabling analytics does not limit the FAS or physical protection of access points.\\\bottomrule
\end{tabularx}

Design acceptance targets are: identified ACS/IDS events at SOC within 5 s on a working network; seven days of recording at the capped bitrate; 30 min of UPS operation for IT/ACS; 24 h standby followed by 30 min FAS alarm, or 24 h followed by 15 min IDS alarm. These are project targets, not requirements quoted from unavailable standards. A single LAN star and one SOC computer remain single points of failure; no reconvergence or server redundancy is claimed.
